\structelem{ЗАКЛЮЧЕНИЕ}

В результате выполнения выпускной квалификационной работы по теме
<<Методика обнаружения несанкционированных действий и аномальной
активности в корпоративной сети с использованием нейронных сетей>>
получены следующие основные результаты.

Все девять задач исходной постановки ВКР выполнены. Соответствие
<<задача --- результат>>:

\begin{enumerate}
    \item Проведен аналитический обзор предметной области
    обнаружения аномалий и несанкционированных действий в
    корпоративных сетях; выполнен сравнительный анализ 9 нейросетевых
    архитектур и 10 публично доступных датасетов NIDS по балльной
    шкале.

    \item Формализованы функциональные и нефункциональные требования
    к системе обнаружения (10 FR + 10 NFR) и определены
    количественные критерии качества с числовыми порогами.

    \item Спроектированы методика обнаружения аномальной активности
    из 10 шагов и архитектура нейросетевого детектора CNN-LSTM
    (Conv1D~$\times$~2 + BiLSTM + Attention pool + LayerNorm).

    \item Разработана подсистема предобработки и формирования
    обучающих и тестовых выборок (log1p + robust scaling,
    one-hot~/~frequency-encoding, скользящие окна $W = 32$, $S = 8$).

    \item Реализована и обучена нейросетевая модель детекции в едином
    программном регистре с 15 моделями-кандидатами; целевая
    CNN-LSTM достигла F1\,=\,0,9735 и PR-AUC\,=\,0,9945.

    \item Разработан конечно-автоматный агент активного тестирования
    (11 состояний) с библиотекой из 7 параметризованных шаблонов и
    инжектором дрейфа трех типов (covariate~/~prior~/~concept).

    \item Реализован программный прототип системы обнаружения в
    реальном времени с REST-интерфейсом (FastAPI) и визуальной
    панелью (Streamlit).

    \item Проведена экспериментальная проверка эффективности методики
    (6 экспериментов E1--E6 и cross-evaluation на CICIDS2017) и
    подтверждена гипотеза исследования; FPR в 11 раз ниже XGBoost при
    сопоставимом F1.

    \item Выполнена оценка эксплуатационных характеристик
    (latency 1,16~мс/окно, throughput 24\,560~EPS на CPU) и оформлены
    итоговые материалы.
\end{enumerate}

Цель работы достигнута. Разработана методика нейросетевого
обнаружения несанкционированных действий и аномальной активности
в корпоративной сети, валидируемая конечно-автоматным агентом
активного тестирования с библиотекой параметризованных шаблонов
атак. Методика подкреплена воспроизводимым программным прототипом
и экспериментально проверена на реальном датасете UNSW-NB15 с
cross-evaluation на CICIDS2017.

По сравнению с лидером среди классических подходов
(XGBoost, F1\,=\,0,965) целевая модель обеспечивает в
одиннадцать раз меньшее число ложных срабатываний на типовом
потоке корпоративного SOC за счет существенно более низкого FPR
на операционном пороге (0,0095 против 0,108). Это формирует
фундаментальный практический аргумент в пользу CNN-LSTM как
proposed-модели для производственного развертывания в составе
подсистемы NTA-аналитики корпоративного SOC.

\medskip

\noindent{Разработанная методика обнаружения аномальной
активности и несанкционированных действий в корпоративной сети.}
По итогам работы сформирована методика, представляющая собой
последовательность взаимосвязанных процедур и формализованных
правил, обеспечивающих воспроизводимое обнаружение аномальной
активности в потоковом сетевом трафике. Методика состоит из десяти
шагов.

\begin{enumerate}
    \item Сбор и валидация входной телеметрии. На вход системы
    поступают потоковые записи (flow) в формате CSV, NetFlow~v9
    или IPFIX. Каждая запись валидируется по Pydantic-схеме на
    соответствие типам и диапазонам признаков; некорректные строки
    отбрасываются с логированием.
    \item Очистка и нормализация. Применяются: унификация регистра
    и пробелов в категориальных полях; устранение дубликатов;
    обрезка хвостов распределений по квантилям 0,001 и 0,999;
    лог-преобразование тяжелохвостых числовых признаков; робастное
    масштабирование (медиана, IQR). Все статистики собираются
    только на обучающей выборке.
    \item Кодирование категориальных признаков. Применяется
    one-hot кодирование для признаков малой кардинальности
    (\texttt{proto}, \texttt{state}) и frequency-кодирование для
    \texttt{service}. Неизвестные категории получают значение
    \texttt{unknown} с нулевой частотой.
    \item Формирование скользящих окон. Из последовательности
    flow-записей формируются окна длины $W = 32$ с шагом $S = 8$
    (75\,\% перекрытие); метка окна определяется по последней
    записи (last-агрегация).
    \item Нейросетевой инференс. Каждое окно подается на вход
    обученной модели CNN-LSTM (двунаправленный LSTM,
    attention-пулинг, LayerNorm), которая возвращает вероятность
    атаки $p \in [0, 1]$ для всего окна.
    \item Калибровка вероятностей и порога. Применяется
    температурное масштабирование логитов с параметром $T$,
    подобранным на валидационной выборке через L-BFGS;
    операционный порог $\tau$ определяется бинарным поиском по
    отсортированному списку скорингов нормальных окон под целевую
    частоту ложных срабатываний $\operatorname{FPR}^* = 0{,}01$.
    \item Принятие решения. Для каждого окна выносится бинарное
    решение $\hat{y} = \mathbb{I}[p \geq \tau]$.
    \item Формирование алертов. Каждое положительное решение
    обрабатывается алертформером, который присваивает уровень
    критичности по пятиступенчатой шкале (info, low, medium, high,
    critical) на основе нормированного превышения скоринга над
    порогом, и выполняет временную дедупликацию повторов по
    триплету (src\_ip, dst\_ip, attack\_cat) в окне 5\,с.
    \item Мониторинг дрейфа распределений. Параллельно с инференсом
    подсистема \texttt{drift\_report} вычисляет индексы PSI, KL и
    MMD$^2$ между текущим окном и эталонной обучающей выборкой;
    при превышении порога $\operatorname{PSI} > 0{,}25$
    устанавливается индикатор \texttt{drift\_alarm}, направляемый
    в SOC как событие дрейфа.
    \item Активное тестирование. На стенде проводится
    воспроизводимое стресс-тестирование детектора с использованием
    конечно-автоматного агента, имеющего 11 состояний (NORMAL,
    7 ATTACK, 3 DRIFT), стохастическую матрицу переходов,
    минимальные длительности состояний, библиотеку из семи
    параметризованных шаблонов атак и инжектор дрейфа трех типов.
    По прогону рассчитываются per-state recall, F1 и FPR.
    \item Экспорт и интеграция. Алерты выгружаются в JSON Lines
    формате, совместимом с типовыми SIEM, а также доступны через
    REST API (\texttt{GET /alerts/recent}) и визуальную панель
    Streamlit, привязанную к сервису по \texttt{DIPLOMA\_API\_URL}.
\end{enumerate}

Сформированная методика обеспечивает повторяемость результатов
обнаружения на одних и тех же входных данных (побитовая
воспроизводимость при фиксированных seed-ах), независимость от
конкретного источника телеметрии (поддерживаются NetFlow~v9 и
IPFIX без модификации кода), прозрачную процедуру калибровки порога
под локальный эксплуатационный FPR и встроенный механизм
сигнализации о дрейфе распределений.
